\section{Retained data and fixed-panel scope}\label{sec:data}
The primary experiment uses retained silver v2 mandi reports, combining the AGMARKNET Kaggle archive before October 2025 with official data.gov.in reports thereafter. Prices are INR per quintal. Input Parquet bytes and shard lists are hashed before and after feature execution; these retained bytes, rather than an assurance of identical later API responses, define the reproduced dataset. Official resource access requires a key, which is excluded from source and archives.

The 36-commodity panel is an existing table of 11,845 eligible market--commodity--variety series, supplied as a fixed, hashed external input. Its original selection is not reproduced or independently certified here. Depending on horizon and cleaning regime, the evaluation contains 4,602--5,182 series, rather than all panel members. Findings apply to these eligible report-days and cannot represent every Indian market or farmer.
Both cleaning variants are newly built from the same inputs, with history beginning 2022-10-01, for warm-up before 2024 training. The cleaned variant removes reports flagged by a fivefold deviation from the prior ten-report geometric mean; the unfiltered variant retains them. Other validity checks, harmonisation, panel and build settings are held fixed. This does not create manual error labels: unusual prices can be genuine. Common origins may still carry different targets and features after filtering, so the comparison is pipeline sensitivity, not a causal effect of cleaning.

Build records bind the committed source, executed assembly JAR, commands, runtime versions and exit codes to output checksums. This is observed execution attestation, not an independent build witness. Large inputs are retained locally outside Git; a public immutable archival package and independent-environment reproduction remain outstanding. The earlier national-scale source audit and full-size experiments are historical supporting material in the historical supplement, not acceptance-critical validation of this primary study.
